\section*{Hoofdstuk \ref{ch:b2:reaction-enthalpy} --- Reactie-enthalpieën}

\begin{solution}{exo:b2:reaction-enthalpy:1}
$\Delta_r H^\circ = 3(-393.51) + 4(-285.83) - (-104.7) - 5(0) =
\qty{-2219.2}{kJ/mol}$, negatief: \omterm{def:b2:reaction-enthalpy:exothermic}{exotherm}. (De gemeten verbrandingsenthalpie
is dezelfde, \qty{-2219.2}{kJ/mol}.)
\end{solution}

\begin{solution}{exo:b2:reaction-enthalpy:2}
Methaan: $890.3/16.04 = \qty{55.5}{kJ/g}$, dus \qty{55.5}{MJ} per
kilogram. Diwaterstof: $285.83/2.016 = \qty{141.8}{kJ/g}$,
\qty{141.8}{MJ/kg}: ongeveer 2.6 keer meer per kilogram (maar veel minder per
liter, want het gas is licht).
\end{solution}

\begin{solution}{exo:b2:reaction-enthalpy:3}
$\Delta\nu_{\text{gas}} = 6 - 7.5 = -1.5$, dus $\Delta_r H = \Delta_r U +
\Delta\nu_{\text{gas}}RT = -3264 - 1.5 \times 8.314 \times 298.15 \times
10^{-3} = \qty{-3267.7}{kJ/mol}$, in overeenstemming met de waarde berekend uit
vormingsenthalpieën, \qty{-3267.6}{kJ/mol}.
\end{solution}

\begin{solution}{exo:b2:reaction-enthalpy:4}
De doelreactie is de eerste reactie min tweemaal de tweede:
$\Delta_r H^\circ = -393.5 - 2(-283.0) = \qty{+172.5}{kJ/mol}$. Ze is
\omterm{def:b2:reaction-enthalpy:exothermic}{endotherm}: de hete cokes koelt af terwijl hij koolstofdioxide omzet in
koolstofmonoxide.
\end{solution}

\begin{solution}{exo:b2:reaction-enthalpy:5}
Verbroken bindingen: C--H en Cl--Cl, $416 + 243 = \qty{659}{kJ/mol}$; gevormde
bindingen: C--Cl en H--Cl, $350 + 432 = \qty{782}{kJ/mol}$: schatting
\qty{-123}{kJ/mol}. Uit vormingsenthalpieën: $-83.7 - 92.3 + 74.9 =
\qty{-101.1}{kJ/mol}$. Het verschil van \qty{22}{kJ/mol} komt van de gemiddelde
C--Cl-waarde, een gemiddelde over verscheidene moleculen, en van de C--H-binding die in
methaan verbroken wordt, sterker dan het gemiddelde van de vier.
\end{solution}

\begin{solution}{exo:b2:reaction-enthalpy:6}
Bij \qty{298}{K}: $-241.83 + 285.83 = \qty{44.00}{kJ/mol}$. Kirchhoff met
$\Delta C_p = 33.59 - 75.35 = \qty{-41.76}{J/(K.mol)}$:
$44.00 - 0.04176 \times 101.85 = \qty{39.75}{kJ/mol}$ bij \qty{400}{K}. De
tabellen geven $-242.85 + 282.59 = \qty{39.74}{kJ/mol}$: de constante
warmtecapaciteiten zijn over dit interval uitstekend.
\end{solution}

\begin{solution}{exo:b2:reaction-enthalpy:7}
Ionisatie: $4.341 \times 96.485 = \qty{418.8}{kJ/mol}$; elektronenopname:
$-3.613 \times 96.485 = \qty{-348.6}{kJ/mol}$. \omterm{def:b2:reaction-enthalpy:lattice-enthalpy}{Roosterenthalpie} $= 436.7 +
89.0 + 121.3 + 418.8 - 348.6 = \qty{717}{kJ/mol}$, minder dan de
\qty{787}{kJ/mol} van natriumchloride: \ce{K+} is groter dan \ce{Na+}, de
ionen liggen verder uit elkaar en trekken elkaar minder aan.
\end{solution}

\begin{solution}{exo:b2:reaction-enthalpy:8}
Warmtecapaciteit van de calorimeter: $C = 2.00/0.418 = \qty{4.785}{kJ/K}$.
Vrijgekomen warmte: $4.785 \times 0.583 = \qty{2.789}{kJ}$, voor $\xi =
\qty{0.0500}{mol}$ (het zuur is beperkend, de base in lichte overmaat):
$\Delta_r H = -2.789/0.0500 = \qty{-55.8}{kJ/mol}$.
\end{solution}

\begin{solution}{exo:b2:reaction-enthalpy:9}
\ce{C8H18(l) + 25/2O2(g) -> 8CO2(g) + 9H2O(l)}: $\Delta_r H^\circ =
8(-393.51) + 9(-285.83) + 250.3 = \qty{-5470.3}{kJ/mol}$; met $M =
\qty{114.2}{g/mol}$ is dat \qty{47.9}{MJ/kg}. Koolstofdioxide: $8 \times 44.01 =
\qty{352.1}{g}$ voor \qty{5.470}{MJ}, dus \qty{64.4}{g/MJ}, tegen
$44.01/0.8903 = \qty{49.4}{g/MJ}$ voor methaan: methaan heeft meer waterstof per
koolstofatoom, en de verbranding van waterstof levert warmte zonder koolstofdioxide.
\end{solution}

\begin{solution}{exo:b2:reaction-enthalpy:10}
Producten van \ce{CH4 + 2O2 -> CO2 + 2H2O(g)}: $C_p = 37.13 + 2 \times 33.59 =
\qty{104.3}{J/K}$; $q = \qty{802.3}{kJ}$; $T_f = 298 + 802\,300/104.3 \approx
\qty{7990}{K}$, bijna driemaal de waarde in lucht, omdat er geen stikstof is
die zijn deel van de warmte opneemt. De echte zuurstofvlam is veel kouder: lang vóór \qty{7990}{K} dissociëren koolstofdioxide en water
tot koolstofmonoxide, waterstof, zuurstof en atomen, \omterm{def:b2:reaction-enthalpy:exothermic}{endotherme}
reacties die het grootste deel van de warmte opnemen; bovendien nemen de warmtecapaciteiten toe met
$T$.
\end{solution}

\begin{solution}{exo:b2:reaction-enthalpy:11}
$\Delta_r H^\circ(700) = \Delta_r H^\circ(298) + \int_{298}^{700}(\Delta_r a +
\Delta_r b\,T)\,\dd T = -91.88 + [-60.5 \times 401.85 + 0.0270 \times
(700^2 - 298.15^2)] \times 10^{-3} = -91.88 - 24.31 + 10.83 =
\qty{-105.4}{kJ/mol}$, op \qty{0.2}{kJ/mol} na gelijk aan de tabellen
(\qty{-105.2}{kJ/mol}), terwijl de schatting met constante $C_p$ er
\qty{4.5}{kJ/mol} naast zat.
\end{solution}

\begin{solution}{exo:b2:reaction-enthalpy:12}
$\Delta_r H^\circ = 2 \times 90.29 = \qty{+180.6}{kJ/mol}$ (\omterm{def:b2:reaction-enthalpy:exothermic}{endotherm}).
$\Delta_r C_p^\circ = 2(29.85) - 29.12 - 29.38 = \qty{1.2}{J/(K.mol)}$, dus
tussen \qty{298}{K} en \qty{2000}{K} verandert de \omterm{def:b2:reaction-enthalpy:reaction-quantity}{reactie-enthalpie} met
ongeveer $1.2 \times 1702 = \qty{2}{kJ/mol}$, één procent. Een \omterm{def:b2:reaction-enthalpy:exothermic}{endotherme}
reactie wordt begunstigd door een hoge temperatuur (\cref{ch:b2:equilibrium-shifts}):
koude lucht maakt bijna geen stikstofmonoxide, de heetste vlammen maken er het
meeste; daarom is de verbrandingstemperatuur een hefboom tegen deze
verontreiniging.
\end{solution}

\begin{solution}{pb:b2:reaction-enthalpy:1}
\textbf{1.} \ce{C4H10(g) + 13/2O2(g) -> 4CO2(g) + 5H2O(l)}.
\textbf{2.} $\Delta_r H^\circ = 4(-393.51) + 5(-285.83) + 125.6 =
\qty{-2877.6}{kJ/mol}$.
\textbf{3.} Met waterdamp: $4(-393.51) + 5(-241.83) + 125.6 =
\qty{-2657.6}{kJ/mol}$; het verschil, \qty{220.0}{kJ}, is de verdampingsenthalpie
van vijf mol water (elk \qty{44.0}{kJ/mol}).
\textbf{4.} $2877.6/58.12 = \qty{49.5}{kJ/g}$ en $2657.6/58.12 =
\qty{45.7}{kJ/g}$.
\textbf{5.} De tweede: het water vertrekt als damp en zijn condensatiewarmte
gaat verloren.
\textbf{6.} $230 \times 45.7 = \qty{1.05e4}{kJ}$, ongeveer \qty{10.5}{MJ}.
\textbf{7.} Bij constant volume is de opgenomen warmte $\Delta U$ (geen
volumearbeid): de bom meet $\Delta_r U$.
\textbf{8.} $\Delta\nu_{\text{gas}} = 4 - 1 - 6.5 = -3.5$ (het water is
vloeibaar).
\textbf{9.} $\Delta_r U = \Delta_r H - \Delta\nu_{\text{gas}}RT = -2877.6 +
3.5 \times 2.479 = \qty{-2868.9}{kJ/mol}$.
\textbf{10.} $\xi = 0.500/58.12 = \qty{8.60e-3}{mol}$; warmte $8.60 \times
10^{-3} \times 2868.9 = \qty{24.68}{kJ}$; $\Delta T = 24.68/10.00 =
\qty{2.468}{K}$.
\textbf{11.} Warmte $10.00 \times 2.461 = \qty{24.61}{kJ}$: $\Delta_r U =
-24.61/(8.603 \times 10^{-3}) = \qty{-2860.6}{kJ/mol}$ en $\Delta_r H =
-2860.6 - 8.7 = \qty{-2869.3}{kJ/mol}$, in absolute waarde \qty{0.29}{\percent} onder de
tabellen: warmteverlies van de calorimeter of een licht
onvolledige verbranding.
\textbf{12.} $n = 1000/18.02 = \qty{55.5}{mol}$; $Q = 55.5 \times 75.35
\times 80 = \qty{334.5}{kJ}$.
\textbf{13.} Vrijgekomen warmte $16.0 \times 45.72 = \qty{731.5}{kJ}$; rendement
$334.5/731.5 = 0.46$.
\textbf{14.} Naar de hete verbrandingsgassen die langs de pan ontsnappen, de
door de vlam verwarmde lucht en de wanden van de pan, en naar straling.
\textbf{15.} $230 \times 45.72 \times 0.457/334.5 = \qty{14.4}{L}$.
\textbf{16.} Met \qty{6.5}{mol} \ce{O2} komen $6.5 \times 78.08/20.95 =
\qty{24.23}{mol}$ \ce{N2} en $6.5 \times 0.93/20.95 = \qty{0.289}{mol}$
\ce{Ar} mee.
\textbf{17.} \qty{4}{mol} \ce{CO2}, \qty{5}{mol} \ce{H2O(g)},
\qty{24.23}{mol} \ce{N2}, \qty{0.289}{mol} \ce{Ar}.
\textbf{18.} De vlam is adiabatisch en isobaar, dus $\Delta H = 0$. Langs een
weg bestaande uit de reactie bij \qty{298}{K} ($\Delta H_1 =
\qty{-2657.6}{kJ}$) gevolgd door het opwarmen van deze producten is $\Delta H_2 =
-\Delta H_1 = \qty{+2657.6}{kJ}$: alle warmte gaat naar de gassen.
\textbf{19.} $C_p = 4(37.13) + 5(33.59) + 24.23(29.12) + 0.289(20.79) =
\qty{1028}{J/K}$.
\textbf{20.} $T_f = 298 + 2\,657\,600/1028 = \qty{2883}{K}$.
\textbf{21.} $T_f = 2300 + 100 \times (2657.6 - 2515.7)/(2655.7 - 2515.7) =
\qty{2401}{K}$.
\textbf{22.} De warmtecapaciteiten van de gassen nemen toe met de temperatuur (met
ongeveer een kwart tussen \qty{298}{K} en \qty{2400}{K}): met de waarden bij
kamertemperatuur worden de gassen te goedkoop opgewarmd.
\textbf{23.} Boven \qty{2000}{K} dissociëren koolstofdioxide en water gedeeltelijk
(\omterm{def:b2:reaction-enthalpy:exothermic}{endotherm}), de vlam straalt, en de menging met lucht is nooit volmaakt.
\textbf{24.} De \omterm{def:b2:reaction-enthalpy:flame-temperature}{adiabatische vlamtemperatuur} van butaan in lucht is
$\boldsymbol{T_{\text{ad}} \approx \qty{2.40e3}{K}}$.
\end{solution}
